\section{Введение}

На прошлой лекции мы закончили говорить о взвешенных графах.
Далее мы поговорим о нахождении LCA -- наибольший общий родитель.
Внезапно это имеет достаточно много применений, в частности, это нам поможет в комбинации с алгоритмом Фараха-Колтона-Бендера ускорить решение задачи static RMQ.

\section{LCA}

\Def Пусть дано подвешененное дерево $T \langle V, E \rangle$. Пусть $u, v \in V$.
Тогда $\texttt{LCA(}u, v\texttt{)}$ будем называть наиболее глубокую вершину такую, что она лежит одновременно на путях от $u$ и $v$ к корню дерева.

Решим задачу нахождения \texttt{LCA} с помощью \textit{метода двоичных подъемов}. Далее всюду $T\langle V, E\rangle$ -- некоторое подвешенное за вершину $r$ дерево.
Глубину будем обозначать 
